% ============================================================
% sections/model_solution.tex — 模型求解 / Model Solution
% ============================================================
% Describe how each model is solved:
%   - Algorithm used / 使用的算法
%   - Software / 软件（MATLAB, Python, Lingo, etc.）
%   - Solution steps / 求解步骤
%   - Results / 求解结果
%   - Convergence proof (if applicable) / 收敛性证明

\section{模型求解}

\subsection{问题一模型的求解}

\subsubsection{求解方法}

[说明采用的求解算法及选择原因。例如：由于模型为线性规划模型，采用单纯形法求解，使用 Lingo/MATLAB 软件实现。]

\subsubsection{求解步骤}

\begin{enumerate}
    \item [步骤1描述]
    \item [步骤2描述]
    \item [步骤3描述]
\end{enumerate}

\subsubsection{求解结果}

[呈现求解结果。用表格或文字描述。]

\begin{table}[htbp]
\centering
\caption{问题一求解结果}
\label{tab:result1}
\begin{tabular}{ccc}
\toprule
\textbf{[列1标题]} & \textbf{[列2标题]} & \textbf{[列3标题]} \\
\midrule
[值1] & [值2] & [值3] \\
[值1] & [值2] & [值3] \\
\bottomrule
\end{tabular}
\end{table}

\subsection{问题二模型的求解}

\subsubsection{求解方法}

[说明求解方法]

\subsubsection{求解步骤}

[列出求解步骤]

\subsubsection{求解结果}

[呈现求解结果。可在正文中放关键结果表，详细结果放附录。]

% Figure placeholder example / 图片占位示例
% \begin{figure}[htbp]
%     \centering
%     \includegraphics[width=0.8\textwidth]{figures/figure_name.png}
%     \caption{[图题]}
%     \label{fig:label}
% \end{figure}

\subsection{问题三模型的求解}

\subsubsection{求解方法}

[说明求解方法]

\subsubsection{求解步骤}

[列出求解步骤]

\subsubsection{求解结果}

[呈现求解结果]
